\chap{A complete regularization diagnostic for the escaping model}\label{sec:diagonal-regularization}
The abstract alternatives become explicit for the diagonal model. They also answer a natural objection to the hard prefix corrections: perhaps a smoother correction could fit increasingly well without escaping. It can fit more efficiently, but it cannot avoid the obstruction.

\subsection{The entire regularized family is explicit}
Use $Tu=(u_n/n)$ and $r=w=(1/n)$ as in Section~\ref{sec:escape}. Choose $\alpha=m^{-2}$ for an integer $m\ge1$. The normal equation in each coordinate gives
\[
 a^{(m)}_n=\frac{m^2}{m^2+n^2},\qquad
 v^{(m)}_n=r_n-(Ta^{(m)})_n=\frac{n}{m^2+n^2}.
\]
The repaired ambient candidate is $(-a^{(m)},1)$. Its protected coordinate is one for every $m$.

\begin{proposition}[Energy, residual, and certified tails]\label{WR10}
The displayed $a^{(m)}$ belongs to $\ell^2$ and is the unique regularized correction at $\alpha=m^{-2}$. Set
\[
 A_m=\norm{a^{(m)}}^2=\sum_{n\ge1}\frac{m^4}{(m^2+n^2)^2},\qquad
 B_m=\norm{v^{(m)}}^2=\sum_{n\ge1}\frac{n^2}{(m^2+n^2)^2}.
\]
Then
\begin{equation}\label{eq:soft-bounds}
 \frac m4\le A_m\le\frac{4m}{3},\qquad
 \frac1{4(m+1)}\le B_m\le\frac2m.
\end{equation}
For the sums $A_{m,K},B_{m,K}$ through any integer $K\ge m$,
\begin{equation}\label{eq:soft-tails}
 0\le A_m-A_{m,K}\le\frac{m^4}{3K^3},\qquad
 0\le B_m-B_{m,K}\le\frac1K.
\end{equation}
\end{proposition}
\begin{proof}
For $n>m$, one has $(a_n^{(m)})^2\le m^4/n^4$, so the correction sequence is square-summable. The diagonal operator is bounded with norm one, and $w=(1/n)\in\ell^2$, hence the residual sequence is also in $\ell^2$. The operator is self-adjoint, and the coordinate equality
\[
 \frac1n\frac{n}{m^2+n^2}
 =m^{-2}\frac{m^2}{m^2+n^2}
\]
proves $T^*v^{(m)}=\alpha a^{(m)}$ as an equality of Hilbert vectors. Theorem~\ref{WR06} therefore identifies this actual infinite vector as the unique minimizer, rather than just a coordinatewise candidate.

For $1\le n\le m$, $m^2/(m^2+n^2)\ge1/2$. There are $m$ such coordinates, so $A_m\ge m/4$. For the upper bound, the first $m$ squares are at most one, and the decreasing function $x^{-4}$ gives
\[
 \sum_{n>m}\frac1{n^4}\le\int_m^\infty x^{-4}\,dx=\frac1{3m^3}.
\]
Thus $A_m\le m+m^4/(3m^3)=4m/3$.

For $n\le m$, the residual square is at most $n^2/m^4$. The sum-of-squares estimate
\[
 \sum_{n=1}^m n^2=\frac{m(m+1)(2m+1)}6\le m^3
\]
holds for $m\ge1$, so this part contributes at most $1/m$. For $n>m$, each residual square is at most $1/n^2$, whose tail is at most $1/m$. This proves $B_m\le2/m$.

For the lower bound, $n>m$ implies $m^2+n^2\le2n^2$, hence
\[
 \frac{n^2}{(m^2+n^2)^2}\ge\frac1{4n^2}.
\]
The other direction of the decreasing-function integral comparison gives
\[
 \sum_{n>m}\frac1{n^2}\ge\int_{m+1}^\infty x^{-2}\,dx=\frac1{m+1}.
\]
This proves $B_m\ge1/(4(m+1))$.

Finally, for $K\ge m$, both omitted tails consist of nonnegative terms. Bound the first by $m^4\sum_{n>K}n^{-4}$ and the second by $\sum_{n>K}n^{-2}$. The same upper integral comparisons give the two certified remainders in \eqref{eq:soft-tails}. No finite truncation is substituted for an equality of infinite energies.
\end{proof}
\written{WR10}

Thus the correction norm grows on the order of $\sqrt m$, the residual norm decays on the order of $m^{-1/2}$, and the penalty $\alpha A_m=A_m/m^2$ tends to zero. The smooth correction avoids the hard discontinuity at a prefix boundary, but it still approaches the forbidden constant sequence coordinate by coordinate.

\fixedplate{35_regularized_escape}{At each fixed coordinate, the soft correction approaches one as $m$ tends to infinity. Its transition moves outward as $m$ increases. The chart compares three exact coordinate formulas $m^2/(m^2+n^2)$, not finite approximations to an unspecified optimization procedure. Every displayed sequence has finite energy, but the family's energies are unbounded.}{fig:regularized-escape}

\subsection{The sharp asymptotic correction cost}
The preceding bounds establish divergence and decay. The leading constant can also be computed, which permits a meaningful comparison with hard prefix repair.

\begin{theorem}[A sharp budget asymptotic for the diagonal datum]\label{WR11}
As $m\to\infty$,
\[
 \frac{A_m}{m}\longrightarrow\frac\pi4,\qquad
 mB_m\longrightarrow\frac\pi4.
\]
For the exact budget profile of this infinite problem,
\begin{equation}\label{eq:sharp-budget}
 \lim_{M\to\infty}M\beta_{T,w}(M)=\frac\pi4.
\end{equation}
The hard prefix correction at budget $M=\sqrt N$ instead has
\[
 \sqrt N\,\norm{w-T\mathbf1_N}\longrightarrow1.
\]
\end{theorem}
\begin{proof}
\emph{Step 1: control the infinite Riemann sums uniformly.} Define
\[
 f(x)=\frac1{(1+x^2)^2},\qquad
 g(x)=\frac{x^2}{(1+x^2)^2}.
\]
Then $A_m/m=m^{-1}\sum_{n\ge1}f(n/m)$ and $mB_m=m^{-1}\sum_{n\ge1}g(n/m)$. Fix $L\ge2$ and put $N=\lfloor mL\rfloor$. On $[0,L]$, both functions are uniformly continuous. If $\omega_f(h)$ is a modulus of continuity there, comparison on the $N$ intervals of length $1/m$ gives
\[
 \left|\frac1m\sum_{n=1}^N f(n/m)-\int_0^{N/m}f(x)\,dx\right|
 \le L\omega_f(1/m).
\]
The missing interval $[N/m,L]$ has length below $1/m$ and a bounded integrand, so it contributes a term tending to zero. The same argument applies to $g$. This proves convergence of each truncated sum to its integral on $[0,L]$.

The tails require a bound independent of sufficiently fine discretization. Since $N\ge mL/2$, the inequalities $f(x)\le x^{-4}$ and $g(x)\le x^{-2}$ give
\[
 \frac1m\sum_{n>N}f(n/m)
 \le m^3\sum_{n>N}n^{-4}\le\frac{m^3}{3N^3}
 \le\frac8{3L^3},
\]
\[
 \frac1m\sum_{n>N}g(n/m)
 \le m\sum_{n>N}n^{-2}\le\frac mN\le\frac2L.
\]
The integral tails are at most $1/(3L^3)$ and $1/L$. Given an error tolerance, choose $L$ first to control both discrete and integral tails, then choose $m$ to control the compact-interval sums. This justifies passing from finite-interval Riemann convergence to the full half-line sums. No uncontrolled interchange of two infinite limits is involved.

\emph{Step 2: evaluate the integrals.} The substitution $x=\tan\theta$ gives $dx=\sec^2\theta\,d\theta$ and $1+x^2=\sec^2\theta$, so
\[
 \int_0^\infty f(x)\,dx=\int_0^{\pi/2}\cos^2\theta\,d\theta=\frac\pi4,
 \qquad
 \int_0^\infty g(x)\,dx=\int_0^{\pi/2}\sin^2\theta\,d\theta=\frac\pi4.
\]
Step 1 now proves the two stated asymptotics.

\emph{Step 3: move from regularization parameters to arbitrary budgets.} Set $M_m=\sqrt{A_m}$. Every positive coordinate $m^2/(m^2+n^2)$ increases strictly with $m$, so $A_m$, and hence $M_m$, increases strictly. Since $A_m/m\to c=\pi/4>0$, these budgets tend to infinity and
\[
 \left(\frac{M_{m+1}}{M_m}\right)^2
 =\frac{A_{m+1}/(m+1)}{A_m/m}\frac{m+1}{m}\longrightarrow1.
\]
Proposition~\ref{WR07} identifies $\beta(M_m)=\sqrt{B_m}$. Thus
\[
 M_m\beta(M_m)=\sqrt{(A_m/m)(mB_m)}\longrightarrow c.
\]
For every sufficiently large $M$, choose $m$ with $M_m\le M\le M_{m+1}$. Monotonicity of $\beta$ gives
\[
 M_m\beta(M_{m+1})\le M\beta(M)\le M_{m+1}\beta(M_m).
\]
The lower bound is $(M_m/M_{m+1})[M_{m+1}\beta(M_{m+1})]$ and the upper is $(M_{m+1}/M_m)[M_m\beta(M_m)]$. Both tend to $c$, proving the full budget limit, not merely convergence on a selected subsequence of budgets.

\emph{Step 4: compare the hard-prefix method.} That correction has norm $\sqrt N$ and residual square $\sum_{n>N}n^{-2}$. Integral comparison yields
\[
 \frac{N}{N+1}\le N\sum_{n>N}\frac1{n^2}\le1.
\]
Taking the limit and then a square root gives its leading product one. The two constants are consequently compared at large actual correction budgets, rather than at arbitrarily equated parameter labels.
\end{proof}
\written{WR11}

At large matched correction budgets, the optimal soft repair therefore has about $\pi/4$ times the hard-prefix residual. This improves the efficiency of approximation, not the existence result. Both residuals tend to zero only as the budget tends to infinity. The limiting constant is a calculation for this declared model, not a universal regularization efficiency factor.

\fixedplate{36_budget_product}{The plotted intervals enclose the actual infinite quantities $\sqrt{A_mB_m}$, using finite sums plus the proved tails. The corresponding hard-prefix intervals enclose $\sqrt{N\sum_{n>N}n^{-2}}$ with $N=m$. These have different natural budget sequences. Their limits give the matched-budget asymptotic constants $\pi/4$ and one.}{fig:budget-product}

\subsection{The same operator has an attainable datum}
Change only the residual to $b_n=1/n^2$. Its minimum norm preimage is $a^\dagger_n=1/n\in\ell^2$. At $\alpha=m^{-2}$, the regularized correction is
\[
 c^{(m)}_n=\frac{m^2}{n(m^2+n^2)},\qquad
 a^\dagger_n-c^{(m)}_n=\frac{n}{m^2+n^2}.
\]
The bias square is exactly the previously computed $B_m$. This links the two cases without conflating their existence hypotheses: the escaping residual of one problem is the regularization bias of an attainable problem.

\begin{corollary}[An explicit noisy attainable repair]\label{WR12}
For the attainable datum $b=(1/n^2)$, suppose $\norm{\widetilde b-b}\le\delta$ with $0<\delta\le1$. Choose
\[
 m=\lceil\delta^{-2/3}\rceil,\qquad\alpha=m^{-2}.
\]
Then the noisy regularized correction satisfies
\[
 \norm{a_\alpha(\widetilde b)-a^\dagger}
 \le\sqrt{\frac2m}+\frac{\delta m}{2}
 \le(\sqrt2+1)\delta^{1/3}.
\]
\end{corollary}
\begin{proof}
The exact-data bias has coordinates $n/(m^2+n^2)$, so its squared norm is precisely $B_m$. By Proposition~\ref{WR10},
\[
 \norm{a_{m^{-2}}(b)-a^\dagger}\le\sqrt{2/m}.
\]
Theorem~\ref{WR09} at the same parameter gives
\[
 \norm{a_{m^{-2}}(\widetilde b)-a_{m^{-2}}(b)}
 \le\frac{\delta}{2\sqrt{m^{-2}}}=\frac{\delta m}{2}.
\]
The triangle inequality proves the first bound.

Because $0<\delta\le1$, the number $u=\delta^{-2/3}$ is at least one. The ceiling satisfies $u\le m<u+1\le2u$. Therefore
\[
 \sqrt{2/m}\le\sqrt2\,\delta^{1/3},
 \qquad \frac{\delta m}{2}\le\delta u=\delta^{1/3}.
\]
Adding these estimates proves the displayed rate with every ceiling factor accounted for. The positive-noise condition is needed to define this particular finite integer choice. The statement does not prescribe an infinite $m$ at $\delta=0$, nor does it use a probabilistic assumption on the perturbation.
\end{proof}
\written{WR12}

This is a deterministic guarantee for the specified datum and Hilbert noise bound. It is not a statistical confidence statement. The identical parameter rule in the escaping problem cannot produce convergence to an exact protected Hilbert solution, because no such solution exists.

\subsection{What the calculation permits a machine to certify}
The finite program can check a normal equation, exact preservation, and a quadratic optimality identity. It can return a dual witness ruling out a smaller residual at a given budget. In the infinite diagonal example it can return partial sums together with rational tail enclosures. Those are distinct certificates with distinct conclusions.

None of them alone establishes unrestricted semantic realizability. The bounded-budget theorem needs a genuinely uniform bound. The exactification theorem needs control of candidate movement. The regularization limit needs actual attainment of the projected datum. The repeated appearance of a small residual does not replace any of those hypotheses.

The useful next diagnostic is therefore not just ``Did the residual decrease?'' It is ``What happened to the correction norm, what quantity was preserved, and what existence theorem applies to this datum?''
